\chapter{Fazit und Ausblick}\label{ch:fazit}

\section{Beantwortung der Forschungsfragen}

\subsection*{F1 (Komposition)}

Die elf Permutations-Achsen (Page, Node, Traversal, ValueHandle,
Concurrency, Allocator, Prefetch, Telemetry, ISA, Layout, Reclamation)
spannen einen Konfigurations-Raum von $\sim$5,5 Milliarden Permutationen
auf. Stand-der-Technik-Implementierungen lassen sich durch
algorithm\_profile-XMLs eindeutig rekonstruieren: 30~SOTA-Profile
(8~Tier-1 + 22~Tier-2/3) belegen die Tragfaehigkeit der Achsen-
Konzeption fuer P01--P32.

Die hybride API-Spezialisierung (1~Param $\Rightarrow$
\texttt{std::vector}, 2~Params $\Rightarrow$ \texttt{std::map},
$N>2 \Rightarrow$ \texttt{std::map<K, std::tuple>}) erlaubt die
nahtlose Integration sowohl algorithmus-orientierter
(\texttt{std::map::find}/\texttt{insert}/\texttt{erase}) als auch
container-orientierter (\texttt{std::vector::push\_back}/\texttt{at})
Such-Patterns auf der gleichen Engine-Klasse.

\subsection*{F2 (Mess-Methodik)}

Die Mess-Methodik gliedert sich in zwei Modi:

\begin{itemize}
    \item \textbf{Production-Mode} (\texttt{COMDARE\_EXPERIMENT\_MODE=OFF},
          Default): kein Mess-Overhead, ResultAggregator nicht in der
          ExecutionEngine.
    \item \textbf{Experiment-Mode} (\texttt{COMDARE\_EXPERIMENT\_MODE=ON},
          vom messung\_driver gesetzt): ResultAggregator als
          integraler ExecutionEngine-Member, sammelt pro
          run\_workload-Aufruf die Hardware-Counter.
\end{itemize}

Profile-aware Workload-Routing (V11.2) waehlt pro Permutations-Modul
den Workload-Typ aus dem \texttt{traversal}-Tag des Profils, was
faire Vergleiche zwischen Range-Query- und Point-Lookup-orientierten
Algorithmen sicherstellt.

\subsection*{F3 (Pruefling-Vergleich)}

Der konkrete Vergleich von PRT-ART gegen die 30~SOTA-Profile steht
nach den ersten messung\_driver-Laeufen aus (siehe Ausblick~\S\ref{sec:ausblick}).
Die Architektur ist vollstaendig vorbereitet:

\begin{itemize}
    \item PRT-ART-Profil in
          \texttt{prt\_art/algorithm\_profiles/prtart\_pruefling.profile.xml}
    \item Drei Pflicht-Messreihen-Templates in
          \texttt{Code/experiment\_config/messreihen.xml}
    \item ExperimentDriver mit Auto-Pickup von 30~SOTA-Profilen
          (V10.5)
    \item PrtArtSearchEngineAdapter mit vollstaendiger ABI-Konformitaet
          inkl. notify\_*-Hooks (V11.1) und allen std::map-Operationen
          (V12.2 + V13.5 merge/extract)
\end{itemize}

\section{Limitierungen}

\begin{itemize}
    \item \textbf{Module-Bodies derzeit Stub:} Die Codegen-Pipeline
          generiert ABI-konforme Skelett-Module pro Profil. Die
          eigentliche Such-Algorithmus-Implementation pro Permutation
          bleibt fuer eine Folge-Phase (V15+) offen.
    \item \textbf{Hardware-Validierung ausstehend:} cmake configure ist
          gruen auf MSVC i7-1270P, aber der vollstaendige ctest-Lauf +
          messung\_driver-E2E auf BlockAO-Plattformen (X86\_AVX2/AVX512,
          ARM\_Neon/SVE) steht aus.
    \item \textbf{Manuskript-Auswertungs-Kapitel:} Aktuell methodisch
          beschrieben; konkrete Mess-Resultate folgen.
\end{itemize}

\section{Zukuenftige Arbeit}\label{sec:ausblick}

\textbf{Kurzfristig (V15+):}

\begin{itemize}
    \item E2E-Lauf der drei Pflicht-Messreihen auf realer Hardware
          (V15.2 in der Sprint-Plannung)
    \item Generierung der Mess-Diagramme via diagram\_generator und
          Integration in Kapitel~\ref{ch:auswertung}
    \item Vervollstaendigung der Module-Bodies pro Permutation
\end{itemize}

\textbf{Mittel- bis langfristig:}

\begin{itemize}
    \item Erweiterung um GPU-Adapter (Grace Hopper Cluster, ZIH TUD)
    \item PIM-Allokator-Integration (A16 Lee et al. HPCA 2026)
    \item Formale Verifikation einzelner Permutations-Klassen
          (StarMalloc-Style A13)
\end{itemize}

Die Drei-Repo-Architektur (cache-engine + prt-art + Diplomarbeit-Code)
ermoeglicht es zukuenftigen Diplomanden, weitere Pruefling-Algorithmen
analog zu PRT-ART einzubringen, ohne die Werkzeug-Bibliothek
(cache-engine) zu modifizieren --- ein wichtiges Designprinzip fuer
die Skalierbarkeit der Forschungs-Infrastruktur.
